Longitudinal data raises a recurring question about whether a population hides distinct paths over time and which entities follow each path. Trajectory clustering recovers the paths but returns only cluster membership, while subgroup discovery returns readable rules but usually targets a single scalar rather than a shape traced over a shared calendar. We introduce SEAL, a method for interpretable subgroup discovery on longitudinal data. SEAL groups entities by how their target trajectory accumulates over time and describes every group by a compact conjunctive rule over engineered descriptors, searching for the grouping and its rule jointly so that it reports a group only when a simple rule captures it. We demonstrate SEAL on five longitudinal panels, covering carbon dioxide emissions, pandemic mortality, national income, air quality, and the share of income that Indian households spend on health. In each of them the discovered rules select groups that correspond to accepted structure, such as income and fuel mix for emissions and population age structure for mortality. Permutation tests confirm that these groups are substantially more coherent than chance groupings and that they separate the target itself. An ablation shows that engineered temporal features, which let a rule describe how a covariate moved over the observation window, sharpen the groups precisely when the trajectory structure turns on that movement, most clearly for national income, where they let a simpler rule carve tighter groups. A non-circularity discipline keeps a rule from restating the definition of a ratio target such as a health-spending share.
